\subsection{Normal Subgroups and Quotient Groups}

\begin{definition}
    Given a group $G$, we say $N$ is a normal subgroup and write
    $N \unlhd G$ if for all $g \in G$, $n \in N$ we have $gng^{-1}\in N \iff gN=Ng \iff gNg^{-1}=N$\\
    In this case we can form the quotient group $G\diagup N = \left\{ gN \mid g \in G \right\}$ with $(xN)(yN)=(xy)N$
\end{definition}

\begin{theorem}
    \label[theorem]{thm:normalquoitentgroups}

    Given a group $G$: $N \unlhd G \iff G\diagup N$ is a group 
\end{theorem}
\begin{proof}
    \((\impliedby)\) Assume $G\diagup N$ is a group

    Let $a,b,c,d \in G$ s.t. $aN=cN$ and $bN=dN$ then $abN=cdN$

    $\implies ca^{-1}=j,bd^{-1}=h,abd^{-1}c^{-1}\in N$

    $ahc^{-1}=ah(ja)^{-1}=aha^{-1}j^{-1}\in N \implies aha^{-1}\in N \implies N \unlhd G$

    $(\implies)$ Assume $N \unlhd G$

    Let $a,b,c,d \in G$ s.t. $aN=cN$ and $bN=dN$ $\implies a\in bN \implies a=bn$ same for $c=dh$ and $dn=jd$ since $dN=Nd$

    \begin{align*}
        (aN)(cN)&=(ac)N\\
            &=(bndh)N\\
            &=(bdjhn)N\\
            &=(bd)N\\
            &=(bN)(dN)
    \end{align*}
\end{proof}

\newpage
\subsubsection{Exercises}
\begin{example}
    Show if $[G:H]=2$ then $H \unlhd G$. What is the isom. class of $G\diagup H$?

    Assume $[G:H]=2 \implies gH \sqcup H = G$ for $g \notin H$

    Let $x \in G$

    Case 1: $x \in H \implies xH=H=Hx$

    Case 2: $x \in gH \implies x=gh$ for $h \in H \implies xH=ghH=gH$

    Consider $Hx$ since $g \notin H \implies Hg=gH \implies Hx=H(gh)=g(Hh)=gH=Hx$

    Therefore $H \unlhd G$ since $|G\diagup H|=2$ it is isomorphic to $\Z_2$
\end{example}